\scanpage{19}
% slanted
\margin{\uline{intercalaire}}
\uline{Proposition 3} (Sous réserve qu'on sache -- ou suppose -- que $\sigma'$ est une réflexion \deleted{d'\unsure{axe}} miroir $D_{\sigma'}\ne D_\sigma\unsure{\}}$ -- ce qui ne fait rien autre que la condition unimodulaire!) :

$1^\circ)$ Supposons $D_\sigma\cap D_{\sigma'}\ne\varnothing$, soit $c$ l'unique pt de $D_\sigma\cap D_{\sigma'}$, \deleted{\unsure{à $\varphi$|Alors}} Alors \deleted{\unsure{pour un anneau}} $c$ est l'unique pt fixe de $E$ sous $\Gamma$. Si on prend $c$ comme origine -- de sorte que $E$ devient un vectoriel, $\Gamma$ opérant linéairement --. Le groupe \inserted{\text{$G$}} des autom de $(E,\Gamma,\varphi)$ (i.e. le commutant de $\Gamma$ dans $\mathrm{Aff}(E)$ ou $\mathrm{Gl}(E)$) est égal au groupe des \deleted{\unsure{scalair}} homothéties non nulles, i.e.
\[
k^*\simeq G
\]
-- donc \deleted{$G$ opère \unsure{de façon simplement}} \inserted{\text{$D_\sigma^*$ est un $G$-torseur.}}

\deleted{\unsure{1}}\ $2^\circ)$ Supposons $D_\sigma\parallel D_{\sigma'}$. Soit $T\subset V={}$espace des translations le \deleted{\unsure{droite|droite ?}} \inserted{\text{1-vectoriel}} direction commune de $D_\sigma$ et $D_{\sigma'}$, et identifions $T$ \inserted{\text{$(\simeq k^{\unsure{*|+}})$}} \deleted{au} \inserted{\text{un}} sous-groupe de $\mathrm{Aff}(E)$ (ss-groupe du groupe des transl\unsure{tions|t\textsuperscript{ons}|ations}).
% archive sheet number: 18
